\section{Discussion}
\label{sec:discussion}

\subsection{What the package gives a practitioner}
\label{sec:disc-summary}

\pkg{amore} closes the full relational-event analysis loop inside a single
\proglang{R} package. From a raw event stream a practitioner can: canonicalize
it into an \code{amore\_event\_log}; draw matched non-events under several risk-set
scopes; compute a catalogue of 74 numeric endogenous statistics spanning the
timing and closure variant axes; fit the resulting case-control table
through one \code{rem()} front-end; and, for the linear fit, run the
martingale-residual goodness-of-fit family of \citet{boschi2025gof}. The same
package also simulates: an exact Gillespie scheme and an approximate
$\tau$-leap scheme share one endogenous-state machine, and the simulator emits a
first-class case-control table that feeds \code{rem()} directly.

The expository spine is one likelihood with three predictors. The \code{clogit}
and \code{nn} backends maximize the \emph{same} risk-set softmax
conditional-logistic partial likelihood --- a linear predictor in the first
case, a hand-derived multilayer perceptron in the second --- while the
\code{gam} backend fits the algebraically related degenerate (case-1-control)
binomial logistic regression of \citet{boschi2025rhem} via \pkg{mgcv}, giving
access to smooth, time-varying, and random-effect terms. The validation battery
(Section~\ref{sec:validation}) supports the implementation: linear effects are
recovered with bias $\le 0.006$ for $|\beta| \le 0.4$ and Wald coverage between
0.90 and 1.00 (E1); a non-linear smooth is recovered with RMSE 0.075 and
correlation 0.998 (E2); the simulator and the post-hoc feature engine agree to
$4.3\times 10^{-14}$ under history-log alignment (E5); and the neural gradients
match numerical differentiation to $1.4\times 10^{-10}$ (E6). The contribution
is the integrated, reproducible artifact and its unifying frame --- not new
statistical theory.

\subsection{Limitations}
\label{sec:disc-limitations}

We state the boundaries of the package plainly, because several of them
constrain how the tool should be used.

\paragraph{Goodness of fit is linear-only.} The four \code{gof\_*} functions and
\code{martingale\_residuals()} operate on the linear case-control fit and hard-reject
anything else; there is no goodness-of-fit machinery for the smooth (\code{gam})
or neural (\code{nn}) backends. The recentering step that aligns the residual
process is a heuristic, and we have not conducted a calibration or power study of
the tests themselves.

\paragraph{The neural backend is deliberately minimal.} The \code{nn} backend is
pure R and GPU-free, training full-batch by default with optional
mini-batching. Its analytic gradients are verified against
numerical differentiation, but the design carries real costs: it returns no
coefficients, standard errors, or any other uncertainty quantification
(\code{coef()} is \code{NULL}); it rejects factor covariates and accepts numeric
inputs only; and the pure-R training loop sets a practical scalability ceiling below a
compiled engine --- mitigated, but not removed, by the mini-batch option. The interaction-recovery
advantage reported in E6 (concordance 0.81 for \code{nn} versus 0.52 for
\code{clogit} on a planted $x_1 x_2$ truth, against ties of 0.78/0.78 on a linear
truth) is an \emph{in-sample} metric on a single fixture; it is evidence that the
joint MLP can represent structure the additive backends cannot, not a
generalization guarantee. We frame this backend strictly as an accessibility and
engineering contribution, with no priority claim: STREAM
\citep{filippimazzola2024} already fits flexible additive effects on the same
case-control objective at far larger scale, and neural temporal-point-process
models \citep{du2016,mei2017} long predate this package.

\paragraph{The simulator has small-network failure modes.} The scaling study (E4)
shows $\tau$-leap is 22--42$\times$ faster than Gillespie at 60 actors and up to
roughly $100\times$ at 100 actors. But both schemes break down at \emph{small}
actor counts under a strongly self-exciting process, for opposite reasons:
Gillespie's per-dyad softmax underflows (``too few positive probabilities'') and
$\tau$-leap's intensities blow up into a memory overflow. A practitioner
simulating small, strongly self-exciting networks should expect to tune the time
step, reduce the effect strength, or accept that some configurations are not
reachable with either engine.

\paragraph{Naming and release status.} The package name shares letters with the
archived CRAN package \code{AMORE} (a neural-network toolbox), which must be
resolved before a CRAN submission. As of this writing the headline backends and
the full statistics engine are present in the source tree; a clean CRAN
re-release exposing them through the installed binary is a prerequisite for the
package to deliver its advertised interface to a user running
\code{install.packages("amore")}.

\paragraph{Hyper-event support is rem-ready but undemonstrated.} The relational
hyper-event (RHEM) feature path --- hyperedge simulation and hyperedge feature
computation --- exists and produces case-control tables of the shape
\code{rem()} consumes. We have not yet shipped a vignette carrying a single
end-to-end \code{rem()} fit on hyper-event data, so RHEM should be read as a
designed-and-wired capability rather than a demonstrated one.

\subsection{Future work}
\label{sec:disc-future}

Several extensions follow directly from the limitations above and would
strengthen the package without altering its frame.

\begin{itemize}
  \item A compiled, minibatched neural engine (e.g.\ a \pkg{torch} backend)
        would lift the pure-R scalability ceiling and make larger event logs
        tractable for the \code{nn} backend.
  \item Factor support via embeddings would let the neural backend consume the
        categorical actor and dyad attributes it currently rejects.
  \item Uncertainty quantification for the neural fit --- bootstrap, or an
        approximate posterior --- would turn its point predictions into
        inferential statements comparable to the linear backend's Wald intervals.
  \item Extending the goodness-of-fit family beyond the linear case-control fit,
        to the smooth and neural backends, would close the most conspicuous gap
        in the current toolchain.
  \item Penalization and basis selection for the additive-spline
        architecture (currently unpenalized, with fixed degrees of freedom per
        covariate, as in \citealp{filippimazzola2024}), and resampling-based
        uncertainty bands for its fitted curves.
  \item A worked end-to-end \code{rem()} fit on hyper-event data would convert the
        RHEM path from a wired capability into a demonstrated one.
\end{itemize}

In its present form \pkg{amore} consolidates a methodology that previously lived
as scattered scripts and a separate simulation tool into one coherent, tested,
documented package, organized around a single unifying view of the
relational-event likelihood. That consolidation, and the reproducible evidence
that the pieces work together correctly, is the contribution we claim.
